% =====================================================================
% Introduction — the case for exact interaction ground truth.
%
% GPS: transformer interactions are non-additive by construction →
% current methods assume additivity → we provide ground truth and
% benchmark methods against it.
%
% Status: v1 first draft. Compiles with main.tex.
% =====================================================================

\section{Introduction}
\label{sec:introduction}

Circuit discovery in transformers identifies a small subnetwork that
explains a model's behavior on a target task. The dominant methods ---
activation patching~\citep{wang2022ioi}, Edge Attribution Patching
(EAP)~\citep{syed2023attribution}, and Automatic Circuit DisCovery
(ACDC)~\citep{conmy2023acdc} --- rank components by individual
importance and assemble circuits by thresholding those scores. The
resulting circuits are additive by construction: each component
contributes independently, and the circuit's total effect is the sum of
its parts.

Attention mechanisms violate this assumption at the architecture level.
\citet{leemann2024slalom} prove that transformers with attention
mechanisms structurally cannot represent additive models over sequences
of more than one token: the softmax normalization forces every head's
output to depend on every other head in the same layer. The theorem
applies to both single-layer and multi-layer architectures, and the
result is tight --- the non-additivity is a property of the function
class, not a failure of training. Any circuit discovery method that
reports only marginal importance scores discards this interaction
structure by design.

The gap between additive discovery and non-additive computation raises a
concrete question: how much interaction structure do different discovery
methods retain, and does the discarded structure carry mechanism? Answering
the question requires exact ground truth for the interaction between every pair
of components in a circuit --- data that, until this work, has not been
available for any real trained model.

We provide this ground truth by computing exhaustive Walsh--Hadamard
coalition sweeps over subcircuits of GPT-2 small on the indirect object
identification (IOI) task. The Walsh--Hadamard transform decomposes the
$2^n$ coalition function values into a spectrum of interaction
coefficients, where each order-2 coefficient $w_{ij}$ measures how much
the joint effect of heads $i$ and $j$ deviates from the sum of their
individual effects. The decomposition is exact, reference-free, and
requires no regularization or held-out validation --- properties
inherited from the transform's orthogonality.

Three results emerge from benchmarking discovery methods against this
ground truth:

\begin{enumerate}
\item \textbf{Weight geometry predicts nothing.} OV composition, QK
  composition, and layer distance --- the standard weight-space features
  from \citet{elhage2021mathematical} --- carry no cross-validated
  predictive power for pairwise interaction coefficients
  (Section~\ref{sec:weight-geometry}). Interaction is a property of the
  forward pass, not the weight matrices.

\item \textbf{Different methods recover different structure.} EAP edge
  scores agree moderately with the Walsh interaction spectrum (Spearman
  $\rho = 0.58$), capturing the linear component of pairwise
  interaction. Path patching, which isolates direct causal edges through
  the residual stream, measures a nearly orthogonal quantity (signed
  Pearson $r = 0.07$). The two methods impose different selective
  pressures on which interactions survive
  (Section~\ref{sec:path-patching}).

\item \textbf{The interaction structure carries mechanism.} Applying a
  complementation test from genetics to the Walsh interaction matrix
  rediscovers that negative name movers invert their contribution sign
  when their upstream S-inhibition heads are removed --- a dependency
  that \citet{wang2022ioi} observed as unexplained
  (Section~\ref{sec:what-epistasis-contains}).
\end{enumerate}

Compressed sensing makes the ground truth practical to compute beyond
small circuits: LASSO recovers the full pairwise interaction spectrum of
a 20-head circuit from 140 random coalitions ($r > 0.95$), a
$7{,}490\times$ compression over exhaustive enumeration
(Section~\ref{sec:sparse-walsh-recovery}). For 30-head circuits, the
cost is roughly 500 forward passes --- a ten-minute GPU job rather than
an impossible $2^{30}$ sweep.

The selective-pressure framing applies beyond the methods tested here.
Any circuit discovery method that reports component-level scores can be
evaluated against the Walsh ground truth by asking how much of the
pairwise interaction spectrum it preserves. The interaction spectrum
itself is a fixed property of the circuit and task, independent of
how it was discovered.
